% GENERATED FILE. Edit canonical lessons, then run npm run generate:books.
% canonical-source-hash: fnv1a64:0df6cd5a297f727a
% canonical-lessons: ES-C513-redaccion, ES-C513-titular, ES-C513-prensa, ES-C513-corresponsal, ES-C513-publicar

\chapter{En el periódico}
\label{ch:es-en-el-periodico}
\hlchaptermodality{513}

\begin{chapteropening}
\textbf{By the end of this chapter:} \emph{I can read a headline and talk about the press, the newsroom, a correspondent and what was published.}

Complete the last lesson of chapter 513: I can read a headline and talk about the press, the newsroom, a correspondent and what was published.
\end{chapteropening}

\section[la redacción]{la redacción — the newsroom; the writing, the essay}

\label{lesson:ES-C513-redaccion}



\begin{quote}
Before the new one: say the Spanish for to come about, then the Spanish for to take place.
\end{quote}

\subsection*{You'll want to know: la redacción}
\textbf{la redacción} — \textquotedblleft{}the newsroom; the writing, the essay\textquotedblright{}.

\begin{quote}
La redacción del periódico está en el centro.
\end{quote}

\begin{quote}
\emph{The paper's newsroom is in the centre.}
\end{quote}

At school \emph{una redacción} is a written essay.

\begin{cousinweb}
\emph{redacción} is built on \emph{redactar} 'to write up, to draw up', which is usually traced to Latin \emph{redāctus}, the participle of \emph{redigere} 'to bring back, to reduce to order' (Lewis \& Short), built on \emph{agere} 'to drive'.
\end{cousinweb}

\begin{grammarlens}[title={What you've built}]
One more word for telling, explaining and describing at B1.
\end{grammarlens}

\subsection*{Your turn}
\begin{itemize}
\raggedright
  \item \emph{Say it:} \emph{la redacción}
  \item \emph{Say it:} \emph{la redacción}, once more
  \item \emph{Say it:} a sentence of your own with \emph{la redacción}
  \item \emph{Recall:} say the Spanish for to come about, then the Spanish for to take place, then say \emph{la redacción} again
\end{itemize}

\begin{tcolorbox}[breakable,colback=teal!4,colframe=teal!35!black,title={Before you move on}]
What does la redacción mean? (\textquotedblleft{}The newsroom; the writing, the essay\textquotedblright{}.) Say it once more.
\end{tcolorbox}
\section[el titular]{el titular — the headline}

\label{lesson:ES-C513-titular}



\begin{quote}
Before the new one: say the Spanish for to take place, then the Spanish for the newsroom.
\end{quote}

\subsection*{You'll want to know: el titular}
\textbf{el titular} — \textquotedblleft{}the headline\textquotedblright{}.

\begin{quote}
El titular habla del accidente.
\end{quote}

\begin{quote}
\emph{The headline is about the accident.}
\end{quote}

\emph{Los titulares} are the headlines. As an adjective, \emph{titular} means holding a post: \emph{el profesor titular}.

\begin{cousinweb}
\emph{titular} comes from Late Latin \emph{titulāre} 'to give a title to' (EtymDB), built on \emph{titulus} 'an inscription, a label' (Lewis \& Short), the word behind \emph{título}. A headline is the title printed over a story.
\end{cousinweb}

\begin{grammarlens}[title={What you've built}]
One more word for telling, explaining and describing at B1.
\end{grammarlens}

\subsection*{Your turn}
\begin{itemize}
\raggedright
  \item \emph{Say it:} \emph{el titular}
  \item \emph{Say it:} \emph{el titular}, once more
  \item \emph{Say it:} a sentence of your own with \emph{el titular}
  \item \emph{Recall:} say the Spanish for to take place, then the Spanish for the newsroom, then say \emph{el titular} again
\end{itemize}

\begin{tcolorbox}[breakable,colback=teal!4,colframe=teal!35!black,title={Before you move on}]
What does el titular mean? (\textquotedblleft{}The headline\textquotedblright{}.) Say it once more.
\end{tcolorbox}
\section[la prensa]{la prensa — the press}

\label{lesson:ES-C513-prensa}



\begin{quote}
Before the new one: say the Spanish for the newsroom, then the Spanish for the headline.
\end{quote}

\subsection*{You'll want to know: la prensa}
\textbf{la prensa} — \textquotedblleft{}the press\textquotedblright{}.

\begin{quote}
La prensa habló mucho del caso.
\end{quote}

\begin{quote}
\emph{The press talked a lot about the case.}
\end{quote}

\emph{La prensa} is the newspapers and the journalists together. \emph{Una rueda de prensa} is a press conference.

\begin{cousinweb}
\emph{prensa} comes from Catalan \emph{premsa} 'a press' (EtymDB), a word usually traced to Latin \emph{premere} 'to press'. The printing press gave its name to the papers it printed.
\end{cousinweb}

\begin{grammarlens}[title={What you've built}]
One more word for telling, explaining and describing at B1.
\end{grammarlens}

\subsection*{Your turn}
\begin{itemize}
\raggedright
  \item \emph{Say it:} \emph{la prensa}
  \item \emph{Say it:} \emph{la prensa}, once more
  \item \emph{Say it:} a sentence of your own with \emph{la prensa}
  \item \emph{Recall:} say the Spanish for the newsroom, then the Spanish for the headline, then say \emph{la prensa} again
\end{itemize}

\begin{tcolorbox}[breakable,colback=teal!4,colframe=teal!35!black,title={Before you move on}]
What does la prensa mean? (\textquotedblleft{}The press\textquotedblright{}.) Say it once more.
\end{tcolorbox}
\section[el corresponsal]{el corresponsal — the correspondent}

\label{lesson:ES-C513-corresponsal}



\begin{quote}
Before the new one: say the Spanish for the headline, then the Spanish for the press.
\end{quote}

\subsection*{You'll want to know: el corresponsal}
\textbf{el corresponsal} — \textquotedblleft{}the correspondent\textquotedblright{}.

\begin{quote}
El corresponsal vive en Roma desde hace años.
\end{quote}

\begin{quote}
\emph{The correspondent has lived in Rome for years.}
\end{quote}

A reporter who sends news from another city or country. The same form for a woman: \emph{la corresponsal}.

\begin{cousinweb}
\emph{corresponsal} is usually traced to \emph{corresponder} 'to correspond, to match', from \emph{co-} and \emph{responder}, which comes from Latin \emph{respondēre} 'to answer', built on \emph{spondēre} 'to promise' (Lewis \& Short).
\end{cousinweb}

\begin{grammarlens}[title={What you've built}]
One more word for telling, explaining and describing at B1.
\end{grammarlens}

\subsection*{Your turn}
\begin{itemize}
\raggedright
  \item \emph{Say it:} \emph{el corresponsal}
  \item \emph{Say it:} \emph{el corresponsal}, once more
  \item \emph{Say it:} a sentence of your own with \emph{el corresponsal}
  \item \emph{Recall:} say the Spanish for the headline, then the Spanish for the press, then say \emph{el corresponsal} again
\end{itemize}

\begin{tcolorbox}[breakable,colback=teal!4,colframe=teal!35!black,title={Before you move on}]
What does el corresponsal mean? (\textquotedblleft{}The correspondent\textquotedblright{}.) Say it once more.
\end{tcolorbox}
\section[publicar]{publicar — to publish}

\label{lesson:ES-C513-publicar}



\begin{quote}
Before the new one: say the Spanish for the press, then the Spanish for the correspondent.
\end{quote}

\subsection*{You'll want to know: publicar}
\textbf{publicar} — \textquotedblleft{}to publish\textquotedblright{}.

\begin{quote}
El periódico publicó la noticia el domingo.
\end{quote}

\begin{quote}
\emph{The newspaper published the news on Sunday.}
\end{quote}

It takes \emph{qu} before \emph{e}: \emph{publiqué}. \emph{Publicar en internet} is to post online.

\begin{cousinweb}
\emph{publicar} is a learned borrowing from Latin \emph{pūblicāre} 'to make public' (EtymDB), built on \emph{pūblicus} 'of the people'. English \emph{publish} comes from the same verb through French.
\end{cousinweb}

\begin{grammarlens}[title={What you've built}]
One more word for telling, explaining and describing at B1.
\end{grammarlens}

\subsection*{Your turn}
\begin{itemize}
\raggedright
  \item \emph{Say it:} \emph{publicar}
  \item \emph{Say it:} \emph{publicar}, once more
  \item \emph{Say it:} a sentence of your own with \emph{publicar}
  \item \emph{Recall:} say the Spanish for the press, then the Spanish for the correspondent, then say \emph{publicar} again
\end{itemize}

\begin{tcolorbox}[breakable,colback=teal!4,colframe=teal!35!black,title={Before you move on}]
What does publicar mean? (\textquotedblleft{}To publish\textquotedblright{}.) Say it once more.
\end{tcolorbox}
